\chapter{Podsumowanie}
\label{chap:summary}

\section*{Cele projektu}
Niniejsza praca miała na celu opracowanie biblioteki,
która wspomaga automatyczny dobór architektury sieci
\gls{pMLP} wraz z~jej parametrami uczenia.
W~ramach realizacji tego stworzona została
biblioteka \acrfull{GAAML} w~języku Python.
Ten cel oraz wszystkie założone szczegółowe
cele zostały zrealizowane.

Została zaprojektowana jednoznaczna reprezentacja genetyczna
struktury sieci wraz z~jej parametrami uczenia.
Zaimplementowane zostały mechanizmy
ewolucyjne oraz ewaluacja modeli.
Dodatkowo biblioteka podczas optymalizacji w~pełni
archiwizuje badane osobniki wraz z~wyuczonymi wagami,
dzięki czemu późniejsza analiza wyników jest dużo łatwiejsza
oraz użytkownik może wybrać, który model sieci chce,
nie tylko ten wyróżniony przez algorytm jako najlepszy.

\section*{Wkład własny}
Wkład własny w~niniejszą pracę, poza niniejszym dokumentem,
wyraża się w~dostarczonym rozwiązaniu praktycznym
w~postaci biblioteki \gls{GAAML}.
Stworzone narzędzie opiera się na autorskim
sposobie krzyżowania oraz funkcji korygującej
uwzględniającej złożoność osobnika.

Biblioteka dzięki temu umożliwia
automatyczną optymalizację,
a~sama implementacja jest modularna i~elastyczna,
co pozwala na stosunkowo prostą integrację
z~istniejącymi projektami czy dalsze jej rozwijanie.

\section*{Streszczenie wyników eksperymentów}
Testy w~\cref{chap:test_and_result}
jednoznacznie potwierdziły tezę pracy.
Algorytm ewolucyjny, operując
w~tym samym i~z~góry określonym budżecie,
wykazał się we wszystkich testach większym
przystosowaniem najlepszego osobnika niż sieci
powstałe z~metody czysto stochastycznej,
czyli przeszukiwania losowego.

\section*{Odpowiedzi na pytania badawcze}
\textbf{%
  Czy zastosowanie algorytmu genetycznego
  pozwala uzyskać modele o~jakości wyższej
  niż modele wybrane przez stochastyczną
  metodę przeszukiwania przestrzeni?%
}\\
Tak, zostało wykazane to przez eksperymenty,
w~których algorytm genetyczny wykazał się
większą skutecznością i~szybszym działaniem.

\textbf{%
  Jak parametry algorytmu genetycznego wpływają
  na skuteczność procesu optymalizacji?%
}\\
Parametry algorytmu takie jak prawdopodobieństwo mutacji
i~krzyżowania definiują jak różnorodne są rozwiązania
oraz jak są rozwiązania eksploatowane.
Współczynnik korekcyjny bezpośrednio wpływa
na to jak ważna jest mała złożoność sieci
w~stosunku do trafnych odpowiedzi modelu.
Liczba osobników w~populacji definiowała
możliwości jednoczesnej eksploracji algorytmu,
a~liczba pokoleń wskazywała jak dużo jest dostępnego
czasu na poszukiwanie optymalnego rozwiązania.

\textbf{%
  Jak zmienia się czas trwania optymalizacji
  w~zależności od rozmiaru populacji i~liczby pokoleń?%
}\\
Niniejsze pytanie nie zostało bezpośrednio
w~niniejszej pracy wykazane,
natomiast podczas dostrajania parametrów
najbardziej oczywistym efektem na długość
uczenia wykazały się liczba osobników w~populacji,
liczba generacji oraz współczynnik korekcyjny.
Pozostałe parametry, czyli współczynnik krzyżowania
oraz mutacji nie miały wyraźnego wpływu na długość procesu.
Zmniejszanie liczby pokoleń bardziej
zmniejszało czas trwania procesu,
natomiast większa liczba populacji dawała większą
bazę startową do szybszej optymalizacji,
dzięki czemu pomimo zwiększania liczby osobników w~populacji
nie zwiększało to proporcjonalnie czasu optymalizacji.

\section*{Ograniczenia projektu}
Poza wielokrotnie wspominanym długim czasem działania,
algorytm z~racji archiwizowania wszystkich wag z~nauczonych
sieci tworzy folder (w~podanej lokalizacji) o~pojemnościach,
które mogą przekraczać $50\text{GB}$.

\section*{Wnioski końcowe i~kierunki dalszego rozwoju}
Podsumowując, przedstawiony projekt
stanowi w~pełni funkcjonalne narzędzie,
które może służyć do dalszych badań i~rozwoju.

Biblioteka \gls{GAAML} może być
rozwijana pod wieloma względami:
\begin{nospitemize}
  \item zaimplementowanie kompresowanie danych,
    by archiwizowane dane nie miały tak dużych objętości,
  \item dodanie obsługiwania architektur rozproszonych,
  \item umożliwienie zatrzymania procesu optymalizacji
    oraz go wznowienia,
  \item umożliwienie wznowienia procesu nawet
    przy natychmiastowym wyłączeniu komputera
    na przykład wskutek wyłączenia prądu.
\end{nospitemize}
